\documentclass[10pt,a4paper]{amsart}
\usepackage{amsmath,amssymb,amsthm,graphicx}
% Retained minimal preamble for the Stacks Project chapter more-algebra.tex (More on Algebra).
% The chapter file pinned-original-source.tex carries only \input{preamble} on its first line, so no
% preamble text lives inside the chapter itself.  The Stacks Project uses one single shared preamble for
% every chapter; the snapshot stages it identically as sources/stacks-chow/preamble.tex and
% sources/stacks-etale-cohomology/preamble.tex.  The block below is copied mechanically, line for line,
% from sources/stacks-chow/preamble.tex lines 14-18 (the Xy-pic setup) and 154-217 (the native theorem
% declarations together with the mathematical macros).  Nothing was added, removed or redefined; the
% wrapper class and the amsmath/amssymb/amsthm/graphicx base packages are supplied by the assembler.
% Chapter: more-algebra.tex -- More on Algebra.
% For commutative diagrams we use Xy-pic
\usepackage[all]{xy}

% We use 2cell for 2-commutative diagrams.
\xyoption{2cell}

% Theorem environments.
%
\theoremstyle{plain}
\newtheorem{theorem}[subsection]{Theorem}
\newtheorem{proposition}[subsection]{Proposition}
\newtheorem{lemma}[subsection]{Lemma}

\theoremstyle{definition}
\newtheorem{definition}[subsection]{Definition}
\newtheorem{example}[subsection]{Example}
\newtheorem{exercise}[subsection]{Exercise}
\newtheorem{situation}[subsection]{Situation}

\theoremstyle{remark}
\newtheorem{remark}[subsection]{Remark}
\newtheorem{remarks}[subsection]{Remarks}

\numberwithin{equation}{subsection}

% Macros
%
\def\lim{\mathop{\mathrm{lim}}\nolimits}
\def\colim{\mathop{\mathrm{colim}}\nolimits}
\def\Spec{\mathop{\mathrm{Spec}}}
\def\Hom{\mathop{\mathrm{Hom}}\nolimits}
\def\Ext{\mathop{\mathrm{Ext}}\nolimits}
\def\SheafHom{\mathop{\mathcal{H}\!\mathit{om}}\nolimits}
\def\SheafExt{\mathop{\mathcal{E}\!\mathit{xt}}\nolimits}
\def\Sch{\mathit{Sch}}
\def\Mor{\mathop{\mathrm{Mor}}\nolimits}
\def\Ob{\mathop{\mathrm{Ob}}\nolimits}
\def\Sh{\mathop{\mathit{Sh}}\nolimits}
\def\NL{\mathop{N\!L}\nolimits}
\def\CH{\mathop{\mathrm{CH}}\nolimits}
\def\proetale{{pro\text{-}\acute{e}tale}}
\def\etale{{\acute{e}tale}}
\def\QCoh{\mathit{QCoh}}
\def\Ker{\mathop{\mathrm{Ker}}}
\def\Im{\mathop{\mathrm{Im}}}
\def\Coker{\mathop{\mathrm{Coker}}}
\def\Coim{\mathop{\mathrm{Coim}}}

% Boxtimes
%
\DeclareMathSymbol{\boxtimes}{\mathbin}{AMSa}{"02}

%
% Macros for moduli stacks/spaces
%
\def\QCohstack{\mathcal{QC}\!\mathit{oh}}
\def\Cohstack{\mathcal{C}\!\mathit{oh}}
\def\Spacesstack{\mathcal{S}\!\mathit{paces}}
\def\Quotfunctor{\mathrm{Quot}}
\def\Hilbfunctor{\mathrm{Hilb}}
\def\Curvesstack{\mathcal{C}\!\mathit{urves}}
\def\Polarizedstack{\mathcal{P}\!\mathit{olarized}}
\def\Complexesstack{\mathcal{C}\!\mathit{omplexes}}
% \Pic is the operator that assigns to X its picard group, usage \Pic(X)
% \Picardstack_{X/B} denotes the Picard stack of X over B
% \Picardfunctor_{X/B} denotes the Picard functor of X over B
\def\Pic{\mathop{\mathrm{Pic}}\nolimits}
\def\Picardstack{\mathcal{P}\!\mathit{ic}}
\def\Picardfunctor{\mathrm{Pic}}
\def\Deformationcategory{\mathcal{D}\!\mathit{ef}}
\title{Switching the order in the determinant det(id + a b) = det(id + b a) for maps between finite projective modules over any ring}
\author{The Stacks Project Authors}
\date{stacks-project:more-algebra.tex:local-snapshot-20261012; GNU FDL 1.2 or later}
\begin{document}
\maketitle
\setcounter{subsection}{1}
\noindent\textit{Editorial note.} All delivered mathematical text is the author's own, copied from the pinned Stacks Project chapter file more-algebra.tex; no statement or proof step was rewritten, reordered or removed. The counted statement is the labeled lemma environment at source lines 36419-36428 and the counted proof is the proof environment at source lines 36430-36455, which immediately follows it and is its own proof.\par\medskip
\noindent\textit{Editorial note.} The lemma printed immediately before this one at source lines 36413-36417 carries a proof whose entire body is the single word Omitted; that neighbouring statement is not part of this card and its span is not counted, so the card contains no omitted proof.\par\medskip
\noindent\textit{Editorial note.} The chapter file pinned-original-source.tex carries only an input of the shared preamble on its first line, so no preamble text lives inside the chapter. The retained file sources/stacks-more-algebra/retained-preamble.tex is a mechanical, line-for-line copy of the project's single shared preamble as staged in the snapshot (Xy-pic setup plus the native theorem declarations and mathematical macros). Nothing was added, removed or redefined, and the excerpt serves only as preamble, never as a span coordinate.\par\medskip
\noindent\textit{Editorial note.} The source URL is the source-level Stacks Project URL, as required when more than one card is produced from one source; the per-claim tag is carried by claim\_key. The counted statement and proof refer to no labelled item of the chapter and cite no bibliography entry, so this card contributes no external reference stub and needs no bibliography.\par\medskip
\noindent\textit{Allowed standard prerequisites (editorial).}\begin{itemize}
\item An equality of elements of a ring of the form det(id + a b) = det(id + b a) may be checked after localizing at every prime, and over a local ring a finite projective module is finite free. used at proof lines 36431-36433, where the author replaces R by a localization at a prime ideal and then assumes R local with M and N finite free. Reduces the statement for arbitrary finite projective modules over an arbitrary ring to the matrix identity of the next step, over a commutative ring of the simplest kind.
\item In characteristic zero the determinant of a square matrix of size at most N is a universal polynomial in the traces of the powers of that matrix. invoked at proof lines 36442-36445, where the author passes to the fraction field of the polynomial ring and then uses that a universal polynomial expresses the determinant in the traces of the powers. Converts the matrix identity for the two rectangular matrices AB and BA into the family of trace identities Trace((I\_n + AB)\textasciicircum{}k) = Trace((I\_m + BA)\textasciicircum{}k), which is what can be expanded.
\item The trace identity Trace(XY) = Trace(YX) for matrices of compatible size, and the factorizations (AB)\textasciicircum{}k = A(BA)\textasciicircum{}\{k-1\}B and (BA)\textasciicircum{}k = (BA)\textasciicircum{}\{k-1\}AB. used at proof lines 36451-36454. Supplies the last step: the traces of all powers agree, hence by the previous item the two determinants agree.
\end{itemize}\medskip
\begin{lemma}
\label{lemma-det-switch}
Let $R$ be a ring. Let $M$, $N$ be finite projective $R$-modules.
Let $a : M \to N$ and $b : N \to M$ be $R$-linear maps.
Then
$$
\det(\text{id} + a \circ b) = \det(\text{id} + b \circ a)
$$
as elements of $R$.
\end{lemma}

\begin{proof}
It suffices to prove the assertion after replacing $R$ by a localization
at a prime ideal. Thus we may assume $R$ is local and $M$ and $N$ are
finite free. In this case we have to prove the equality
$$
\det(I_n + AB) = \det(I_m + BA)
$$
of usual determinants of matrices where $A$ has size $n \times m$
and $B$ has size $m \times n$. This reduces to the case
of the ring $R = \mathbf{Z}[a_{ij}, b_{ji}; 1 \leq i \leq n, 1 \leq j \leq m]$
where $a_{ij}$ and $b_{ij}$ are variables and the entries of the
matrices $A$ and $B$. Taking the fraction field, this reduces to the
case of a field of characteristic zero. In characteristic zero there
is a universal polynomial expressing the determinant of a matrix
of size $\leq N$ in the traces of the powers of said matrix.
Hence it suffices to prove
$$
\text{Trace}((I_n + AB)^k) = \text{Trace}((I_m + BA)^k)
$$
for all $k \geq 1$. Expanding we see that it suffices to prove
$\text{Trace}((AB)^k) = \text{Trace}((BA)^k)$ for all $k \geq 0$.
For $k = 1$ this is the well known fact that
$\text{Trace}(AB) = \text{Trace}(BA)$. For $k > 1$
it follows from this by writing $(AB)^k = A(BA)^{k - 1}B$
and $(BA)^k = (BA)^{k - 1} A B$.
\end{proof}

\end{document}
